\documentclass[11pt]{article}
\usepackage{docstyle}
\renewcommand\sheetname{Mini mock \textperiodcentered{} Green Machines + Metals}
\begin{document}
\sheethead{Mini mock}{Green Machines + Metallic materials \textperiodcentered{} Year 9 Science}{17 minutes \textperiodcentered{} each question: A / M / E}
Answer every question. Use full sentences for every \textbf{explain} question.

\Q Spelling counts in this question.\lvl{A}
\par\sub{a} Name the tubes that carry glucose around a plant. \hfill\blank[40mm]{phloem}
\par\sub{b} Name the sticky part of a flower where pollen lands. \hfill\blank[40mm]{stigma}
\par\sub{c} Name the property of a metal that can be stretched into wires. \hfill\blank[40mm]{ductile}\par\vspace{1mm}

\Q Explain why pea seeds need oxygen to germinate.\lvl{M}
\alines{3}{The seed uses oxygen for respiration. Respiration releases energy from the food stored in the cotyledons, and the seed needs this energy to grow (the radicle and plumule grow out).}

\Q A piece of metal has a mass of 53.4\,g. When it is put into a measuring cylinder, the water rises from 15.0\,mL to 21.0\,mL.\lvl{A / M}
\par\sub{a} Calculate the volume of the metal.
\alines{1}{$21.0 - 15.0 = 6.0$ mL}
\sub{b} Calculate the density of the metal. Give the unit.
\flines{2}{density $= \dfrac{53.4\ \mathrm{g}}{6.0\ \mathrm{mL}} = 8.9$ g/mL (copper)}

\Q Two leaves on the same plant are used. Leaf 1 is covered on both sides with foil. Leaf 2 is not covered. After 2 days in sunlight, both leaves are tested with iodine.\lvl{E}
\par\sub{a} Predict the colour of each leaf after the iodine test.
\alines{1}{Leaf 1 stays brown (orange-brown). Leaf 2 turns blue-black.}
\sub{b} Explain your predictions.
\alines{3}{Leaf 1 gets no light, so it cannot photosynthesise and makes no glucose, so there is no starch. Leaf 2 gets light, photosynthesises and stores the glucose as starch.}
\sub{c} Why was Leaf 2 needed in this experiment?
\alines{2}{It is the control: it shows that a leaf in light does make starch, so the result for Leaf 1 is caused by the missing light.}

\Q Explain why kitchen sinks are made of stainless steel and not ordinary iron.\lvl{M}
\alines{3}{A sink is always wet and in air, so iron would rust (it has water and oxygen). Stainless steel is an alloy of iron with chromium and nickel, which resists corrosion, so it does not rust.}

\Q Tin reacts slowly with dilute hydrochloric acid. Copper does not react with it.\lvl{E}
\par\sub{a} Write the word equation for tin reacting with dilute hydrochloric acid.
\alines{1}{tin + hydrochloric acid $\to$ tin chloride + hydrogen}
\sub{b} Copper is used for water pipes, but iron water pipes need to be protected. Explain why, using the reactivity series.
\alines{3}{Copper is low in the reactivity series, so it does not react with water or air in the pipe. Iron is more reactive: with water and oxygen it rusts, so it must be protected, for example by galvanising.}
\end{document}
